\newpage

\section[METODOLOGIA]{Metodologia}

% Neste capítulo devem ser descritos os procedimentos a serem seguidos na realização da pesquisa. A metodologia depende da natureza do trabalho, do tipo de pesquisa que se pretende desenvolver e, principalmente dos objetivos que se propõem alcançar. De acordo com Gil (2018), na metodologia são apresentadas as seguintes informações: 

% Tipo de pesquisa: esclarecer a natureza (exploratória, descritiva ou explicativa) e o delineamento da pesquisa (experimental, estudo de caso, pesquisa bibliográfica, etc.); 

% População e amostra: apresentar o universo a ser estudado, extensão da amostra e forma de seleção; 

% Coleta de dados: descrever as técnicas a serem utilizadas para coleta de dados (questionários, ensaios laboratoriais, ensaios de campo, técnicas de entrevista ou de observação, etc.); 

% Análise dos dados: descrever os procedimentos a serem adotados para análise quantitativa e/ou qualitativa.

2. FUNDAMENTAÇÃO TEÓRICA
2.1 SOLO - nesse subitem você pode definir o solo (produto do intemperismo sobre as rochas) e fazer comentários sobre sua origem (solos residuais, solos transportados e as características dos horizontes), fazer comentários sobre o tamanho das partículas, comentários sobre o comportamento desse material com a variação da umidade (limites de consistência)
2.2 SOLO COMO MATERIAL CONSTRUTIVO - nesse subitem você pode retomar os benefícios dessas técnicas construtivas, tanto em aspectos ambientais (disponibilidade do material, reciclabilidade, consumo de energia, geração de resíduos) como em relação ao conforto dos usuários (térmico, acústico, etc.). Aqui você pode citar brevemente algumas técnicas como taipa de pilão, cob, pau a pique, bloco de terra comprimida e o próprio adobe. 
2.2.1. Adobe - nesse subitem você pode se aprofundar um pouco mais na técnica do adobe: origem, regiões do planeta em que sua utilização é mais comum, tipo de solo que melhor se adapta à técnica, adições comuns para o melhoramento de suas propriedades (pelos e gordura animal, adição de fibras vegetais, etc) e estratégias para a durabilidade das construções (beirais, impermeabilização das fundações), etc. Nesse subitem foque em práticas tradicionais. O uso de produtos industrializados como cal e cimento serão abordados no próximo subitem.
2.3 TÉCNICAS DE MELHORAMENTO DE SOLOS - nesse subitem você pode contextualizar a necessidade de melhoramento de solos para que ele usado para fins construtivos, você pode trazer os principais tipos de estabilização: mecânica, granulométrica/física e química.
2.3.1 Solo-cimento - nesse subitem você pode trazer mais detalhes sobre o solo-cimento, principais características, tipos de cimento comumente utilizados para esse fim, principais reações químicas, tipos de solo que melhor se adaptam à técnica e principais limitações (tendência de trincamento, perda de resistência quando exposto à umidade excessiva, etc.).